% Pizarra de T3, páginas 1--4.
\begin{ejer}[La sucesión exponencial, a nivel de pregavilla]\label{ejer:t3-sucesion-exponencial}
Sea $X=\mathbb C$ con la topología clásica, $\mathcal O_X$ la gavilla de funciones holomorfas, y $\mathcal F$ la pregavilla de funciones que admiten un logaritmo holomorfo. Describe una sucesión exacta de pregavillas en $X$
\[
0\longrightarrow\mathbb Z\longrightarrow\mathcal O_X\longrightarrow\mathcal F\longrightarrow0,
\]
donde $\mathbb Z\to\mathcal O_X$ es la inclusión natural y $\mathcal O_X\to\mathcal F$ está dada por $f\mapsto\exp(2\pi i f)$. Verifica la exactitud. Muestra que $\mathcal F$ no es una gavilla (falla el axioma de engarce: hay funciones que no tienen logaritmo global pero sí lo tienen localmente). \hfill [Vakil 2.3.J]
\end{ejer}
\begin{sol}
Sea $\mathbb Z$ la gavilla de funciones localmente constantes con valores enteros. Precisaré, también, que
\[
\mathcal F(U)\coloneq\set*{g\colon U\to\mathbb C\given
\exists\widetilde g\en\mathcal O_X(U),\ g=\exp\circ\widetilde g}.
\]
Como toda función localmente constante es holomorfa, la inclusión $I\colon\mathbb Z\hookrightarrow\mathcal O_X$ está bien definida (en cada abierto es la inclusión $\mathbb Z(U)\hookrightarrow\mathcal O_X(U)$; la naturalidad es obvia). Asimismo, si $f\en\mathcal O_X(U)$, entonces $\exp(2\pi i f)$ tiene a $2\pi i f$ como logaritmo holomorfo, por lo que $\exp(2\pi i f)\en\mathcal F(U)$. Luego, $E\colon\mathcal O_X\to\mathcal F$ está bien definida; además, conmuta con restricciones.

Notemos que $\mathbb Z$ y $\mathcal O_X$ son gavillas de grupos abelianos mediante la suma usual, mientras que $\mathcal F$ lo es como pregavilla con la multiplicación. Aclaro esto para notar que en $\mathcal F$ el neutro es $1$, a pesar de la notación $\mathcal F\to0$. Siguiendo esta línea, notemos que cada componente de $I$ es homomorfismo (pues es una inclusión) y también cada componente de $E$ es homomorfismo, ya que
\[
E_U(f+g)=\exp\bigl(2\pi i(f+g)\bigr)
=\exp(2\pi i f)\exp(2\pi i g)=E_U(f)E_U(g).
\]
Esto garantiza que la sucesión corta presentada es de pregavillas de grupos abelianos.

\medskip
\noindent\emph{Probemos exactitud}:\\
Como trabajamos en la categoría de pregavillas, podemos calcular núcleos e imágenes entrada por entrada.
\begin{itemize}
\item Como $I_U$ es inyectiva para cada abierto $U\subseteq X$, tenemos
\[
\forall U\en\Open(X),\qquad (\ker I)(U)=\ker(I_U)=0
\quad\Longrightarrow\quad\ker I=0.
\]
Esto prueba la exactitud en $\mathbb Z$.

\item Dado $U\en\Open(X)$, para todo $g\en\mathcal F(U)$, por definición existe $f_1\en\mathcal O_X(U)$ con $g=\exp\circ f_1$. Luego, $f_2\coloneq\frac{1}{2\pi i}f_1$ cumple que
\[
E_U(f_2)=\exp\left(2\pi i\frac{1}{2\pi i}f_1\right)=\exp(f_1)=g.
\]
Esto prueba que $\mathcal F(U)=\im(E_U)=(\im E)(U)$ para todo $U\en\Open(X)$, por lo que $\mathcal F=\im E$ y la sucesión es exacta en $\mathcal F$.

\item Finalmente, dado $U\en\Open(X)$, por resultados de variable compleja y continuidad:
\[
f\en\ker(E_U)
\quad\Longleftrightarrow\quad\exp(2\pi i f)=1
\quad\Longleftrightarrow\quad f(z)\en\mathbb Z\quad\forall z\en U
\quad\Longleftrightarrow\quad f\en\mathbb Z(U).
\]
Esto prueba $(\ker E)(U)=\mathbb Z(U)$ para cada $U\en\Open(X)$, de donde $\ker E=\mathbb Z=\im I$.
\end{itemize}
Estos tres puntos prueban la exactitud de la sucesión corta de pregavillas.

\medskip
\noindent\emph{Probemos que $\mathcal F$ no es gavilla}:\\
Sea $S\subseteq\Open(X)$ la familia de abiertos simplemente conexos de $X$ que no contienen al cero. Entonces, $S$ es cubierta de $\mathbb C^*\coloneq\mathbb C\setminus\set{0}$. Consideremos la familia $\set{\iota_U\colon U\to\mathbb C}_{U\en S}$ con $\iota_U(z)=z$. Es resultado elemental que todo $U\en S$ admite una rama del logaritmo, es decir, $\log_U\en\mathcal O_X(U)$ con $\iota_U=\exp\circ\log_U$. Esto prueba que $\iota_U\en\mathcal F(U)$; además, la familia $\set{\iota_U}_{U\en S}$ es compatible en traslapes.

No obstante, su pegado como función $\iota_{\mathbb C^*}\colon\mathbb C^*\to\mathbb C$ no admite una rama del logaritmo. Se sigue que $\iota_{\mathbb C^*}\noten\mathcal F(\mathbb C^*)$, por lo que $\set{\iota_U}_{U\en S}$ no admite pegado (si lo tuviera, tendría que ser $\iota_{\mathbb C^*}$).
\end{sol}
